% ========================================================================================
% CAPITOLO 1 - RETI COMBINATORIE ELEMENTARI
% ========================================================================================
\chapter{Reti combinatorie elementari}

% ----------------------------------------------------------------------------------------
\section{Esercizio 1: Multiplexer 16:1}
% ----------------------------------------------------------------------------------------

\subsection{Esercizio 1.1}

\begin{traccia}
	Progettare, implementare in VHDL e testare mediante simulazione un multiplexer
	indirizzabile 16:1, utilizzando un approccio di progettazione per composizione a
	partire da multiplexer 4:1.
\end{traccia}

\progetto

Il multiplexer 16:1 seleziona uno fra 16 ingressi \sig{d(15 downto 0)} in base al valore
della selezione \sig{s(3 downto 0)}: $y = d_s$. Seguendo un approccio \emph{per composizione},
il componente è costruito con cinque istanze del multiplexer 4:1 \ent{mux4}
(Appendice~\ref{app:mux4}) disposte su due livelli (Figura~\ref{fig:mux16}):

\begin{itemize}[nosep]
	\item il \textbf{primo livello} contiene quattro mux 4:1, uno per ogni gruppo di quattro
	      ingressi, pilotati dai bit meno significativi \sig{s(1 downto 0)};
	\item il \textbf{secondo livello} contiene un mux 4:1 che sceglie il gruppo, pilotato dai
	      bit più significativi \sig{s(3 downto 2)}.
\end{itemize}

\begin{figure}[H]
	\centering
	\begin{tikzpicture}
		\foreach \i in {0,...,3} {
			\pgfmathtruncatemacro{\lo}{4*\i}
			\pgfmathtruncatemacro{\hi}{4*\i+3}
			\node[mux, minimum width=1.3cm, minimum height=1cm] (m\i) at (0,{-1.6*\i}) {4:1};
			\draw[bus] (-2.6,{-1.6*\i}) node[left, font=\small]{\sig{d(\hi..\lo)}} -- node[larghezza=4]{/} (m\i.west);
		}
		% selezione del primo livello: una linea comune che raggiunge i quattro mux
		\draw[seg] (-0.9,-5.6) node[below, font=\scriptsize]{\sig{s(1..0)}} -- (-0.9,{-1.6*3-0.75});
		\foreach \i in {0,...,3} {
			\draw[seg] (-0.9,{-1.6*\i-0.75}) -- ({-0.02},{-1.6*\i-0.75}) -- (m\i.south);
			\fill (-0.9,{-1.6*\i-0.75}) circle (1.2pt);
		}
		\draw[seg] (-0.9,{-1.6*3-0.75}) -- (-0.9,{-0.75});
		\node[mux, minimum width=5.6cm, minimum height=1.2cm] (mout) at (3.6,-2.4) {4:1};
		\foreach \i in {0,...,3} {
			\draw[seg] (m\i.east) -- (m\i.east -| mout.west);
		}
		\draw[seg] (mout.south) ++(0,-0.6) node[below, font=\scriptsize]{\sig{s(3..2)}} -- (mout.south);
		\draw[seg] (mout.east) -- ++(1,0) node[right]{\sig{y}};
		\node[font=\sffamily\scriptsize, color=black!60] at (0,1.1) {primo livello};
		\node[font=\sffamily\scriptsize, color=black!60] at (3.6,1.1) {secondo livello};
	\end{tikzpicture}
	\caption{Architettura del multiplexer 16:1 ottenuto per composizione di mux 4:1.}
	\label{fig:mux16}
\end{figure}

\begin{porte}[Interfaccia dell'entity \ent{mux16}]
	\porta{d}{in}{16}{ingressi dati}
	\porta{s}{in}{4}{selezione: indice dell'ingresso da portare in uscita}
	\porta{y}{out}{1}{uscita, $y = d_s$}
\end{porte}

\begin{osservazione}[Perché cinque mux 4:1]
	Un mux $2^k$:1 si ottiene da mux 4:1 con $\lceil k/2 \rceil$ livelli. Per $k=4$ servono
	due livelli: $16/4 = 4$ mux al primo e $4/4 = 1$ al secondo, in totale cinque.
\end{osservazione}

\implementazione

L'architettura strutturale replica il primo livello con un costrutto \texttt{for\ldots generate}
(Listato~\ref{lst:mux16}); il codice del componente \ent{mux4} è in appendice.

\vhdlfile{es1/mux16.vhd}{Multiplexer 16:1 strutturale (\texttt{mux16.vhd}).}{lst:mux16}

\begin{attenzione}[Ordine dei bit di selezione]
	Il primo livello deve usare i bit \emph{meno} significativi della selezione. Scambiando
	\sig{s(1 downto 0)} e \sig{s(3 downto 2)} il circuito resta un mux 16:1 valido, ma con
	gli ingressi in ordine permutato: con $s=1$ verrebbe selezionato $d_4$ invece di $d_1$.
\end{attenzione}

\simulazione

Il testbench (Listato~\ref{lst:tb-mux16}) verifica in modo \emph{esaustivo} tutte le 16
selezioni. Per ogni valore di \sig{s} applica due pattern: uno in cui solo l'ingresso
selezionato vale 1 (l'uscita deve valere 1) e il suo complemento (l'uscita deve valere 0).
Così si verifica sia che venga scelto l'ingresso giusto, sia che gli altri non influenzino l'uscita.

\vhdlfile[firstline=18, lastline=43, firstnumber=18]{es1/tb_mux16.vhd}{Processo di stimolo del testbench (estratto di \texttt{tb\_mux16.vhd}).}{lst:tb-mux16}

\screenshot[0.95\linewidth]{es1_sim_mux16.png}{Forme d'onda della simulazione del mux 16:1: per ogni selezione, l'uscita segue l'ingresso selezionato.}{fig:sim-mux16}

\begin{risultato}[Risultato della simulazione]
	Tutti i 32 casi di test sono superati: la console di XSim riporta
	\emph{``Simulazione terminata con 0 errori''}.
\end{risultato}

% ----------------------------------------------------------------------------------------
\subsection{Esercizio 1.2}

\begin{traccia}
	Utilizzando il componente sviluppato al punto precedente, progettare, implementare in
	VHDL e testare mediante simulazione una rete di interconnessione a 16 sorgenti e 4 destinazioni.
\end{traccia}

\progetto

La rete porta la sorgente \sig{d(src)} sulla destinazione \sig{o(dst)}. È composta dal
\ent{mux16} dell'esercizio precedente, che sceglie la sorgente, seguito da un demultiplexer
1:4 che la instrada verso la destinazione (Figura~\ref{fig:interconnect}).

\begin{figure}[H]
	\centering
	\begin{tikzpicture}[node distance=2.2cm]
		\node[blocco] (mux) {MUX 16:1\\\ent{mux16}};
		\node[blocco, right=of mux] (demux) {DEMUX 1:4};
		\draw[bus] ($(mux.west)+(-1.8,0)$) node[left]{\sig{d}} -- node[larghezza=16]{/} (mux.west);
		\draw[seg] (mux) -- (demux);
		\draw[bus] (demux.east) -- node[larghezza=4]{/} ++(1.8,0) node[right]{\sig{o}};
		\draw[bus] ($(mux.south)+(0,-0.9)$) node[below]{\sig{src}} -- node[larghezza=4]{/} (mux.south);
		\draw[bus] ($(demux.south)+(0,-0.9)$) node[below]{\sig{dst}} -- node[larghezza=2]{/} (demux.south);
	\end{tikzpicture}
	\caption{Rete di interconnessione a 16 sorgenti e 4 destinazioni.}
	\label{fig:interconnect}
\end{figure}

\begin{scelta}[Destinazioni non selezionate]\label{sc:mux-dest}
	Le tre destinazioni non selezionate valgono \texttt{0}. La traccia non specifica il
	loro valore; mantenere il valore precedente richiederebbe dei registri e renderebbe la
	rete sequenziale, mentre l'esercizio riguarda le reti combinatorie.
\end{scelta}

\implementazione
\todo{Codice di \ent{demux1\_4} e della rete \ent{interconnect\_16x4} (\texttt{\textbackslash vhdlfile}).}

\simulazione
\todo{Testbench: tutte le coppie (\sig{src}, \sig{dst}) con pattern di dati significativi.}

% ----------------------------------------------------------------------------------------
\subsection{Esercizio 1.3}

\begin{traccia}
	Sintetizzare ed implementare su board il progetto della rete di interconnessione
	sviluppato al punto 1.2, utilizzando gli switch per fornire gli input di selezione e i
	led per visualizzare i 4 bit di uscita. Per quanto riguarda i 16 bit dato in input, essi
	devono essere immessi mediante switch, 8 bit alla volta, sviluppando un'apposita
	``rete di controllo'' per l'acquisizione che utilizzi due bottoni della board per
	caricare rispettivamente la prima e la seconda metà del dato in ingresso.
\end{traccia}

\sintesiboard

\begin{scelta}[Assegnazione di switch, pulsanti e LED]
	\centering\small
	\begin{tabular}{@{}ll@{}}
		\toprule
		\textbf{Risorsa} & \textbf{Funzione} \\ \midrule
		\sig{SW[7:0]}   & metà del dato da caricare \\
		\sig{SW[11:8]}  & sorgente \sig{src} \\
		\sig{SW[13:12]} & destinazione \sig{dst} \\
		\sig{BTNL}      & carica la metà bassa \sig{d(7..0)} \\
		\sig{BTNR}      & carica la metà alta \sig{d(15..8)} \\
		\sig{LED[3:0]}  & uscite \sig{o(3..0)} \\
		\bottomrule
	\end{tabular}
\end{scelta}

La ``rete di controllo'' per l'acquisizione è formata da due registri a 8 bit abilitati
dagli impulsi generati dai pulsanti (Figura~\ref{fig:es13-top}). Ogni pulsante passa per un
\ent{debouncer} (Appendice~\ref{app:debouncer}), che sincronizza il segnale con il clock,
elimina i rimbalzi e produce un impulso lungo un solo ciclo alla pressione.

\begin{figure}[H]
	\centering
	\begin{tikzpicture}[node distance=1.2cm and 1.6cm]
		\node[blocco, minimum width=1.8cm] (dbl) {debouncer};
		\node[blocco, minimum width=1.8cm, below=0.6cm of dbl] (dbr) {debouncer};
		\node[reg, right=of dbl] (rl) {REG 8};
		\node[reg, right=of dbr] (rh) {REG 8};
		\coordinate (regmid) at ($(rl)!0.5!(rh)$);
		\node[blocco, right=2.6cm of regmid, minimum height=2.6cm] (net) {rete di\\interconnessione};
		\draw[seg] ($(dbl.west)+(-1,0)$) node[left]{\sig{BTNL}} -- (dbl);
		\draw[seg] ($(dbr.west)+(-1,0)$) node[left]{\sig{BTNR}} -- (dbr);
		\draw[seg] (dbl) -- node[above, font=\scriptsize]{en} (rl);
		\draw[seg] (dbr) -- node[above, font=\scriptsize]{en} (rh);
		\draw[bus] ($(rl.north)+(0,0.8)$) node[above]{\sig{SW[7:0]}} -- (rl.north);
		\draw[bus] ($(rh.south)+(0,-0.8)$) node[below]{\sig{SW[7:0]}} -- (rh.south);
		\draw[bus] (rl.east) -- node[above, font=\scriptsize]{\sig{d(7..0)}} (rl.east -| net.west);
		\draw[bus] (rh.east) -- node[above, font=\scriptsize]{\sig{d(15..8)}} (rh.east -| net.west);
		\draw[bus] (net.east) -- ++(1.2,0) node[right]{\sig{LED[3:0]}};
		\draw[bus] ($(net.south)+(0,-0.9)$) node[below]{\sig{SW[13:8]}} -- (net.south);
		\begin{scope}[on background layer]
			\node[draw=black!40, dashed, rounded corners, fit=(dbl)(dbr)(rl)(rh), inner sep=0.35cm,
			      label={[font=\sffamily\scriptsize, black!60, anchor=north west]south west:rete di acquisizione}] {};
		\end{scope}
	\end{tikzpicture}
	\caption{Architettura del top-level per la scheda (clock e reset omessi).}
	\label{fig:es13-top}
\end{figure}

\todo{Codice del top-level \ent{top\_es1\_3}.}

\xdcfile{es1/es1_3.xdc}{File dei vincoli per la Nexys A7-100T (\texttt{es1\_3.xdc}).}{lst:es13-xdc}

\begin{attenzione}[Switch 8 e 9]
	Sulla Nexys A7 gli switch \sig{SW[8]} e \sig{SW[9]} appartengono a un banco alimentato a
	\SI{1.8}{\volt} e richiedono \texttt{IOSTANDARD LVCMOS18}, a differenza di tutti gli altri
	pin usati (\texttt{LVCMOS33}).
\end{attenzione}

\screenshot[0.6\linewidth]{es1_board.jpg}{Prova su scheda: la sorgente 5 è instradata verso la destinazione 2.}{fig:es13-board}

% ----------------------------------------------------------------------------------------
\section{Esercizio 2: Sistema ROM+M}
% ----------------------------------------------------------------------------------------

\subsection{Esercizio 2.1}

\begin{traccia}
	Progettare, implementare in VHDL e testare mediante simulazione un sistema S composto da
	una ROM puramente combinatoria di 16 locazioni da 8 bit ciascuna e da una macchina
	combinatoria M che opera come segue: fornito al sistema un indirizzo A di 4 bit, il
	sistema restituisce il valore contenuto nella ROM all'indirizzo A opportunamente
	``trasformato'' attraverso la macchina M. Il comportamento della macchina M è totalmente
	a scelta dello studente, l'unico vincolo è che essa prenda in ingresso 8 bit e ne
	fornisca in uscita 4.
\end{traccia}

\progetto
\todo{Schema a blocchi ROM $\rightarrow$ M e scelta della funzione M (es.\ conteggio degli uni).}

\implementazione
\todo{Codice di ROM, M e sistema S.}

\simulazione
\todo{Testbench esaustivo sui 16 indirizzi.}

\subsection{Esercizio 2.2}

\begin{traccia}
	Sintetizzare ed implementare su board il progetto del sistema ROM+M sviluppato al punto
	2.1, utilizzando gli switch per fornire l'indirizzo della ROM da cui leggere i valori da
	trasformare e i led per visualizzare i 4 bit di uscita.
\end{traccia}

\sintesiboard
\todo{Top-level, file dei vincoli, foto della prova su scheda.}
